
\documentclass[11pt,letterpaper]{article}

% ============================================================
% COSMIC UNIVERSALISM WEBSITE MASTER IMPLEMENTATION MANUAL
% Compile with: pdfLaTeX
% Current stable website integration verified: July 23, 2026
% ============================================================

% ------------------------------------------------------------
% Core packages
% ------------------------------------------------------------
\usepackage[margin=0.78in,headheight=16pt]{geometry}
\usepackage[T1]{fontenc}
\usepackage[utf8]{inputenc}
\usepackage{lmodern}
\usepackage[scaled=0.94]{helvet}
\usepackage{microtype}
\usepackage{parskip}
\usepackage{setspace}
\usepackage{enumitem}
\usepackage{booktabs}
\usepackage{tabularx}
\usepackage{array}
\usepackage{longtable}
\usepackage{xcolor}
\usepackage{graphicx}
\usepackage{fancyhdr}
\usepackage{lastpage}
\usepackage{titlesec}
\usepackage{hyperref}
\usepackage{xurl}
\usepackage{bookmark}
\usepackage{listings}
\usepackage[most]{tcolorbox}
\usepackage{amssymb}
\usepackage{amsmath}
\usepackage{pifont}
\usepackage{needspace}
\usepackage{caption}
\usepackage{tocloft}

% ------------------------------------------------------------
% Document typography
% ------------------------------------------------------------
\renewcommand{\familydefault}{\sfdefault}
\setstretch{1.08}
\setlength{\parindent}{0pt}
\setlength{\parskip}{0.65em}
\raggedbottom
\urlstyle{same}
\setlength{\cftsecnumwidth}{3.3em}
\setlength{\cftsubsecnumwidth}{4.2em}

% ------------------------------------------------------------
% Cosmic Universalism color palette
% ------------------------------------------------------------
\definecolor{CUNavy}{HTML}{050812}
\definecolor{CUNavyTwo}{HTML}{080D1A}
\definecolor{CUSurface}{HTML}{0C1424}
\definecolor{CUGold}{HTML}{DDBF72}
\definecolor{CUGoldLight}{HTML}{FFF8DD}
\definecolor{CUCyan}{HTML}{35CDEB}
\definecolor{CUCyanLight}{HTML}{78E8FF}
\definecolor{CUBlue}{HTML}{668ED1}
\definecolor{CUBlueLight}{HTML}{9EB9EA}
\definecolor{CUAmber}{HTML}{D99A32}
\definecolor{CURed}{HTML}{B84C4C}
\definecolor{CUGreen}{HTML}{2C8A65}
\definecolor{CULight}{HTML}{F3F6FA}
\definecolor{CUBorder}{HTML}{CCD5E0}
\definecolor{CUText}{HTML}{18212E}
\definecolor{CUMuted}{HTML}{5C6878}
\definecolor{CUCodeBackground}{HTML}{F6F8FB}
\definecolor{CUCodeBorder}{HTML}{BFCADA}

% ------------------------------------------------------------
% Project metadata
% ------------------------------------------------------------
\newcommand{\ProjectTitle}{Cosmic Universalism Website}
\newcommand{\DocumentTitle}{Master Implementation and Operations Manual}
\newcommand{\DocumentSubtitle}{Website, Cosmic Breath, CU-Time, CUCII, AI Project Governance, and LaTeX Authoring Standard}
\newcommand{\RepositoryName}{CosmicUniversalismStatement}
\newcommand{\RepositoryURL}{https://github.com/willmaddock/CosmicUniversalismStatement}
\newcommand{\ProtectedBranch}{main}
\newcommand{\StableBranch}{website}
\newcommand{\WebsiteDirectory}{website/}
\newcommand{\CurrentIntegrationTag}{website-v1.1-cosmic-breath}
\newcommand{\CurrentIntegrationCommit}{7ac8716}
\newcommand{\CurrentIntegrationFullCommit}{7ac87167c51aad9ae7a31e17d1b2127fc069e15c}
\newcommand{\CurrentIntegrationTagObject}{0bc09615dc2356a14f384d3c947c7c81e85ced33}
\newcommand{\PriorReleaseTag}{website-v1.3-aurelius-release}
\newcommand{\PriorReleaseCommit}{5a56085}
\newcommand{\FeatureCommit}{5be2efb}
\newcommand{\MainCommit}{8637cad}
\newcommand{\DocumentVersion}{1.1}
\newcommand{\VerificationDate}{July 23, 2026}
\newcommand{\AuthorName}{William Maddock}

% ------------------------------------------------------------
% Hyperlinks and PDF metadata
% ------------------------------------------------------------
\hypersetup{
    colorlinks=true,
    linkcolor=CUBlue,
    urlcolor=CUBlue,
    citecolor=CUBlue,
    pdftitle={Cosmic Universalism Website Master Implementation and Operations Manual},
    pdfauthor={William Maddock},
    pdfsubject={Unified Astro website, Cosmic Breath, CU-Time, CUCII, ChatGPT/Codex governance, release workflow, and LaTeX documentation standard},
    pdfkeywords={Cosmic Universalism, Cosmic Breath, TOM, Astro, TypeScript, Vitest, CU-Time,
                 CUCII, ChatGPT, Codex, Claude, Aurelius, Alexa+, LaTeX,
                 GitHub Pages, project instructions}
}

% ------------------------------------------------------------
% Header and footer
% ------------------------------------------------------------
\pagestyle{fancy}
\fancyhf{}
\fancyhead[L]{\small\textbf{Cosmic Universalism Master Manual}}
\fancyhead[R]{\small\DocumentTitle}
\fancyfoot[L]{\small Version \DocumentVersion}
\fancyfoot[C]{\small \thepage\ of \pageref{LastPage}}
\fancyfoot[R]{\small Baseline: \texttt{Cosmic Breath}}
\renewcommand{\headrulewidth}{0.5pt}
\renewcommand{\footrulewidth}{0.5pt}

% ------------------------------------------------------------
% Section formatting
% ------------------------------------------------------------
\titleformat{\section}
  {\Large\bfseries\color{CUNavy}}
  {\thesection}
  {0.8em}
  {}
  [\vspace{0.2em}\color{CUCyan}\titlerule]

\titleformat{\subsection}
  {\large\bfseries\color{CUNavyTwo}}
  {\thesubsection}
  {0.7em}
  {}

\titleformat{\subsubsection}
  {\normalsize\bfseries\color{CUBlue}}
  {\thesubsubsection}
  {0.6em}
  {}

\titlespacing*{\section}{0pt}{2.1em}{0.9em}
\titlespacing*{\subsection}{0pt}{1.5em}{0.65em}
\titlespacing*{\subsubsection}{0pt}{1.1em}{0.35em}

% ------------------------------------------------------------
% Lists and tables
% ------------------------------------------------------------
\setlist[itemize]{leftmargin=1.6em,itemsep=0.3em,topsep=0.3em}
\setlist[enumerate]{leftmargin=2em,itemsep=0.4em,topsep=0.35em}
\newcommand{\checkbox}{\(\square\)}
\newcommand{\checked}{\ding{51}}
\newcolumntype{Y}{>{\raggedright\arraybackslash}X}
\newcolumntype{L}[1]{>{\raggedright\arraybackslash}p{#1}}

% ------------------------------------------------------------
% General information boxes
% ------------------------------------------------------------
\newtcolorbox{purposebox}{
    enhanced, breakable, colback=CULight, colframe=CUBlue,
    boxrule=0.8pt, arc=2mm, left=3mm, right=3mm, top=2.5mm, bottom=2.5mm,
    title=\textbf{Purpose}, coltitle=white, colbacktitle=CUBlue
}
\newtcolorbox{checkpointbox}{
    enhanced, breakable, colback=CUCyan!6, colframe=CUCyan,
    boxrule=0.8pt, arc=2mm, left=3mm, right=3mm, top=2.5mm, bottom=2.5mm,
    title=\textbf{Checkpoint}, coltitle=CUNavy, colbacktitle=CUCyanLight
}
\newtcolorbox{safetybox}{
    enhanced, breakable, colback=CUAmber!8, colframe=CUAmber,
    boxrule=0.8pt, arc=2mm, left=3mm, right=3mm, top=2.5mm, bottom=2.5mm,
    title=\textbf{Safety Rule}, coltitle=white, colbacktitle=CUAmber
}
\newtcolorbox{warningbox}{
    enhanced, breakable, colback=CURed!6, colframe=CURed,
    boxrule=0.8pt, arc=2mm, left=3mm, right=3mm, top=2.5mm, bottom=2.5mm,
    title=\textbf{Do Not Continue Until Resolved}, coltitle=white, colbacktitle=CURed
}
\newtcolorbox{successbox}{
    enhanced, breakable, colback=CUGreen!7, colframe=CUGreen,
    boxrule=0.8pt, arc=2mm, left=3mm, right=3mm, top=2.5mm, bottom=2.5mm,
    title=\textbf{Successful Result}, coltitle=white, colbacktitle=CUGreen
}
\newtcolorbox{notebox}{
    enhanced, breakable, colback=CUCodeBackground, colframe=CUBorder,
    boxrule=0.7pt, arc=2mm, left=3mm, right=3mm, top=2.5mm, bottom=2.5mm
}
\newtcolorbox{authoritybox}{
    enhanced, breakable, colback=CUGold!8, colframe=CUGold,
    boxrule=0.8pt, arc=2mm, left=3mm, right=3mm, top=2.5mm, bottom=2.5mm,
    title=\textbf{Authority and Source Priority}, coltitle=CUNavy, colbacktitle=CUGoldLight
}
\newtcolorbox{authoringbox}{
    enhanced, breakable, colback=CUBlue!6, colframe=CUBlue,
    boxrule=0.8pt, arc=2mm, left=3mm, right=3mm, top=2.5mm, bottom=2.5mm,
    title=\textbf{Manual Authoring Standard}, coltitle=white, colbacktitle=CUBlue
}

% ------------------------------------------------------------
% Listing styles
% ------------------------------------------------------------
\lstdefinestyle{terminalstyle}{
    basicstyle=\ttfamily\small,
    backgroundcolor=\color{CUNavyTwo},
    keywordstyle=\color{CUCyanLight},
    stringstyle=\color{CUGoldLight},
    commentstyle=\color{CUBlueLight},
    identifierstyle=\color{white},
    showstringspaces=false,
    columns=fullflexible,
    keepspaces=true,
    breaklines=true,
    breakatwhitespace=false,
    frame=none,
    tabsize=2,
    upquote=true
}
\lstdefinestyle{promptstyle}{
    basicstyle=\ttfamily\small,
    backgroundcolor=\color{CUCodeBackground},
    keywordstyle=\color{CUBlue},
    stringstyle=\color{CUGreen},
    commentstyle=\color{CUMuted},
    identifierstyle=\color{CUText},
    showstringspaces=false,
    columns=fullflexible,
    keepspaces=true,
    breaklines=true,
    breakatwhitespace=false,
    frame=none,
    tabsize=2,
    upquote=true
}
\lstdefinestyle{codestyle}{
    basicstyle=\ttfamily\small,
    backgroundcolor=\color{CUCodeBackground},
    keywordstyle=\color{CUBlue},
    stringstyle=\color{CUGreen},
    commentstyle=\color{CUMuted},
    identifierstyle=\color{CUText},
    showstringspaces=false,
    columns=fullflexible,
    keepspaces=true,
    breaklines=true,
    breakatwhitespace=false,
    frame=none,
    tabsize=2,
    upquote=true
}

\newtcblisting{commandbox}[1]{
    enhanced, breakable, listing only, listing engine=listings,
    listing options={style=terminalstyle},
    colback=CUNavyTwo, colframe=CUCyan, coltitle=CUNavy,
    colbacktitle=CUCyanLight, boxrule=0.8pt, arc=2mm,
    left=2mm, right=2mm, top=1mm, bottom=1mm,
    title=\textbf{#1}
}
\newtcblisting{promptbox}[1]{
    enhanced, breakable, listing only, listing engine=listings,
    listing options={style=promptstyle},
    colback=CUCodeBackground, colframe=CUBlue, coltitle=white,
    colbacktitle=CUBlue, boxrule=0.8pt, arc=2mm,
    left=2mm, right=2mm, top=1mm, bottom=1mm,
    title=\textbf{Copy-and-Paste Prompt: #1}
}
\newtcblisting{codebox}[1]{
    enhanced, breakable, listing only, listing engine=listings,
    listing options={style=codestyle},
    colback=CUCodeBackground, colframe=CUCodeBorder, coltitle=white,
    colbacktitle=CUNavyTwo, boxrule=0.8pt, arc=2mm,
    left=2mm, right=2mm, top=1mm, bottom=1mm,
    title=\textbf{#1}
}

% ------------------------------------------------------------
% Convenience commands
% ------------------------------------------------------------
\newcommand{\filepath}[1]{\nolinkurl{#1}}
\newcommand{\branchname}[1]{\texttt{#1}}
\newcommand{\statusline}[2]{\textbf{#1:} #2}
\newcommand{\term}[1]{\textbf{\textsf{#1}}}
\newcommand{\release}[1]{\texttt{#1}}
\newcommand{\CurrentRoute}{/CosmicUniversalismStatement/}

\begin{document}
\hypersetup{pageanchor=false}

% ============================================================
% Cover page
% ============================================================
\begin{titlepage}
\pagecolor{CUNavy}
\color{white}
\thispagestyle{empty}

\vspace*{0.55in}

{\small\bfseries\color{CUCyanLight}
COSMIC UNIVERSALISM COMPUTATIONAL INTELLIGENCE INITIATIVE}

\vspace{0.38in}

{\Huge\bfseries
\ProjectTitle}

\vspace{0.16in}

{\LARGE\color{CUGoldLight}
\DocumentTitle}

\vspace{0.12in}

{\large\color{CUBlueLight}
\DocumentSubtitle}

\vspace{0.34in}

{\large
A unified, controlled manual that consolidates the website foundation, the
Cosmic Breath Cycle Explorer, the Cosmic Breath Time Converter, the CUCII
Prompt Studio, the Aurelius Claude-proven configuration, ChatGPT/Codex project
governance, release operations, and the canonical LaTeX style for future
Cosmic Universalism manuals.}

\vfill

\begin{tcolorbox}[
    colback=CUNavyTwo, colframe=CUBlue, boxrule=0.8pt, arc=2mm,
    width=\textwidth
]
\color{white}
\begin{tabularx}{\textwidth}{>{\bfseries}L{1.75in}Y}
Repository & \RepositoryName \\
Protected branch & \texttt{\ProtectedBranch} \\
Stable website branch & \texttt{\StableBranch} \\
Current integration tag & \texttt{\CurrentIntegrationTag} \\
Current website commit & \texttt{\CurrentIntegrationCommit} \\
Website directory & \texttt{\WebsiteDirectory} \\
Technology & Astro, TypeScript, HTML, modern CSS, Markdown, SVG, lightweight browser JavaScript \\
Verified quality gate & 251 tests across 30 files; Astro static build; 8 pages \\
Document version & \DocumentVersion \\
Verified & \VerificationDate \\
\end{tabularx}
\end{tcolorbox}

\vspace{0.32in}

{\large\bfseries \AuthorName}

\vspace{0.08in}

{\small\url{\RepositoryURL}}

\end{titlepage}
\hypersetup{pageanchor=true}
\pagenumbering{arabic}
\nopagecolor
\color{CUText}

% ============================================================
% Document Control
% ============================================================
\section*{Document Control}
\addcontentsline{toc}{section}{Document Control}

\begin{tabularx}{\textwidth}{>{\bfseries}L{1.95in}Y}
Document & Cosmic Universalism Website: Master Implementation and Operations Manual \\
Version & \DocumentVersion \\
Classification & Master technical specification, operations runbook, AI project bootstrap, and documentation standard \\
Repository & \href{\RepositoryURL}{\RepositoryName} \\
Protected branch & \branchname{\ProtectedBranch} \\
Stable integration branch & \branchname{\StableBranch} \\
Current integration tag & \release{\CurrentIntegrationTag} at \release{\CurrentIntegrationCommit} \\
Prior application release & \release{\PriorReleaseTag} at \release{\PriorReleaseCommit} \\
Current website commit & \release{\CurrentIntegrationFullCommit} \\
Protected main commit & \release{\MainCommit} \\
Website source directory & \filepath{\WebsiteDirectory} \\
Primary local editor & PyCharm \\
Primary implementation environment & ChatGPT Codex / local project chat \\
Version-control interface & PyCharm terminal, Git CLI, or GitHub Desktop \\
Deployment target & GitHub Pages \\
Technical baseline verified & \VerificationDate \\
\end{tabularx}

\begin{purposebox}
This document replaces the need to begin every new project chat with separate
website, Cosmic Breath, CU-Time, and CUCII implementation manuals. It preserves
their essential technical, mathematical, prompt-architecture, testing, and
release rules while adding the completed Cosmic Breath explorer, sealed
structural and content authorities, learning and empirical layers, the current
website integration state, reusable ChatGPT/Codex project-file templates, and
the canonical LaTeX house style for future manuals.
\end{purposebox}

\begin{safetybox}
Earlier manuals remain valuable historical specifications, but their feature
branches, baseline tags, test counts, and ``current'' boundaries may be
superseded. A new chat must use this document's Document Control page and the
live Git repository as the current authority. Never revive an old feature
branch merely because it appears in a historical manual.
\end{safetybox}

\begin{authoritybox}
When sources disagree, use this priority:
\begin{enumerate}
    \item Live repository state verified with Git commands.
    \item Current website integration tag and commit listed in this manual.
    \item This master manual's protected rules and current-state record.
    \item Feature-specific control specimens and reference implementations.
    \item Earlier implementation manuals as historical design records.
    \item Research documents as CU source material, not as silent authorization
          to change production code.
\end{enumerate}
\end{authoritybox}

\clearpage
\tableofcontents
\clearpage

% ============================================================
\section*{START HERE - Initialize a New Project Chat}
\addcontentsline{toc}{section}{START HERE - Initialize a New Project Chat}

\begin{successbox}
A new chat can operate safely from this manual alone for general website
planning. For exact converter parity or exact prompt-fixture work, also provide
the authoritative HTML or text control specimen named by the relevant section.
\end{successbox}

\subsection*{Required reading order}

Before proposing edits, the new chat must read:

\begin{enumerate}
    \item Document Control and Source Priority.
    \item Part 1: Governance and current release.
    \item Part 2: Technology stack and repository architecture.
    \item The feature division relevant to the requested work.
    \item Part 10: Testing and acceptance.
    \item Part 11: Git and release operations.
    \item Part 12 when creating or revising a manual.
\end{enumerate}

\subsection*{Current project identity}

\begin{tabularx}{\textwidth}{>{\bfseries}L{1.85in}Y}
Repository & \RepositoryURL \\
Local repository & \filepath{/Users/cosmos/Documents/GitHub/CosmicUniversalismStatement} \\
Website directory & \filepath{website/} \\
Protected branch & \branchname{main} \\
Stable website branch & \branchname{website} \\
Current integration tag & \release{website-v1.1-cosmic-breath} \\
Current website commit & \filepath{7ac87167c51aad9ae7a31e17d1b2127fc069e15c} \\
Prior application release & \release{website-v1.3-aurelius-release} at \release{5a56085} \\
Current Cosmic Breath route & \filepath{/CosmicUniversalismStatement/cosmic-breath/} \\
Current CUCII route & \filepath{/CosmicUniversalismStatement/cu-intelligence/} \\
Current CU-Time route & \filepath{/CosmicUniversalismStatement/cu-time/} \\
\end{tabularx}

\subsection*{New-chat bootstrap prompt}

\begin{promptbox}{Initialize the Cosmic Universalism Website Project Chat}
We are continuing controlled work on the Cosmic Universalism website.

Read the Master Implementation and Operations Manual before proposing
or making changes. Treat its current release record and governance
rules as mandatory.

Repository:
/Users/cosmos/Documents/GitHub/CosmicUniversalismStatement

Protected branch:
main

Stable website branch:
website

Current website integration tag:
website-v1.1-cosmic-breath

Current website integration commit:
7ac87167c51aad9ae7a31e17d1b2127fc069e15c

Prior application release:
website-v1.3-aurelius-release at 5a56085

Before editing:

1. Confirm the repository root and origin.
2. Confirm the active branch and HEAD.
3. Confirm the working-tree status.
4. Identify the exact feature scope.
5. Inspect only the files relevant to that scope.
6. Report the smallest safe implementation plan.
7. Distinguish current live state from historical manual baselines.

Do not edit files during the first inspection.

Do not reset, clean, stash, discard, pull, switch branches, commit,
push, merge, tag, release, or deploy unless explicitly authorized.

Do not alter CU-Time mathematics or approved source URLs without a
separate documented review.

Do not claim that CUCII retrains, permanently aligns, or changes the
beliefs, policies, memory, or underlying capabilities of an external
AI platform.
\end{promptbox}

\subsection*{Required terminal preflight}

\begin{commandbox}{Verify the live repository before every editing session}
cd /Users/cosmos/Documents/GitHub/CosmicUniversalismStatement || exit 1

git remote get-url origin
git branch --show-current
git rev-parse --short HEAD
git status --short --branch
git tag --points-at HEAD
\end{commandbox}

\begin{warningbox}
Stop when the branch, commit, tag, or working-tree state differs from the
authorized task. Do not use reset, clean, stash, or branch switching to
``repair'' an unexplained difference.
\end{warningbox}

% ============================================================
\section{Part 1 - Governance, Releases, and Source Authority}
% ============================================================

\subsection{Repository policy}

The repository is \href{\RepositoryURL}{\RepositoryName}. The
\branchname{main} branch is protected from website-development work. The
stable website integration branch is \branchname{website}. New development is
performed on a feature branch created from a verified stable website commit.

The normal relationship is:

\begin{codebox}{Controlled branch relationship}
main                         protected research baseline
website                      stable website integration branch
feature/<feature-name>       isolated implementation and testing branch
website-vX.Y-<approved-name> permanent annotated recovery or integration tag
\end{codebox}

\subsection{Current verified release history}

\begin{tabularx}{\textwidth}{L{2.35in}L{1.0in}Y}
\toprule
\textbf{Tag} & \textbf{Commit} & \textbf{Meaning} \\
\midrule
\release{website-v1.1-cu-time-release} & \release{3fa4a37} & Stable CU-Time website release and pre-CUCII recovery point. \\
\release{website-v1.2-cucii-release} & \release{ae16ae7} & Stable CUCII Prompt Studio release. \\
\release{website-v1.3-aurelius-release} & \release{5a56085} & Stable Aurelius Claude preset, Alexa+, and Claude public-demonstration release. \\
\release{website-v1.1-cosmic-breath} & \release{7ac8716} & Current Cosmic Breath integration tag. The numeric label identifies this feature line; it does not imply a downgrade or chronological replacement of the earlier application v1.3 tag. \\
\bottomrule
\end{tabularx}

The current \branchname{website} integration commit is
\filepath{7ac87167c51aad9ae7a31e17d1b2127fc069e15c}. The annotated tag object is
\filepath{0bc09615dc2356a14f384d3c947c7c81e85ced33}, and its peeled commit is the
current website integration commit.

Cosmic Breath feature-history commits include:

\begin{itemize}
    \item \release{d9d4d5c}: canonical 51-state structural ledger.
    \item \release{95b46a7}: typed structural state engine.
    \item \release{38c43f3}: accessible explorer shell.
    \item \release{d6a6ee5}: drag, snap, and settle behavior.
    \item \release{0631b69}: adaptive Canvas atmosphere.
    \item \release{131d792}: companion content authority.
    \item \release{5be2efb}: completed explorer, learning content, empirical context, and navigation.
\end{itemize}

\subsection{Current verification record}

At the July 23, 2026 Cosmic Breath integration checkpoint:

\begin{itemize}
    \item 251 automated tests passed across 30 test files.
    \item The Astro production build completed successfully with no warnings or errors.
    \item Eight static pages were generated, including the Cosmic Breath, CUCII, and CU-Time routes.
    \item The feature branch was merged with an explicit two-parent merge commit.
    \item Local and remote \branchname{website} both resolved to \filepath{7ac87167c51aad9ae7a31e17d1b2127fc069e15c}.
    \item The annotated \release{website-v1.1-cosmic-breath} tag was pushed and verified by its peeled commit.
    \item The protected \branchname{main} branch remained unchanged at \filepath{8637cad8dd208f6e202f74d7fbb74ef5266a2aff}.
    \item The working tree was clean and synchronized; no deployment or GitHub Release occurred.
\end{itemize}

\subsection{Historical manuals and current state}

The source manuals documented controlled development at different times:

\begin{itemize}
    \item the local website foundation;
    \item the CU-Time converter migration;
    \item the CUCII AI Exploration and Prompt Studio;
    \item the Cosmic Breath Interactive Cycle Explorer planning and implementation specification.
\end{itemize}

Their architecture, safety, and workflow rules remain useful. Their historical
feature branch names and expected test counts do not override the current
integration record above. Master Manual v1.0 remains the immutable July 21
historical baseline; this v1.1 source and PDF are the current operational
revision.

\subsection{Phase authorization rule}

Every implementation phase follows four gates:

\begin{enumerate}
    \item \textbf{Inspect}: read live files and report current behavior.
    \item \textbf{Plan}: identify the smallest safe file set and acceptance
          criteria.
    \item \textbf{Implement}: edit only after explicit authorization.
    \item \textbf{Verify}: run tests, build, browser checks, and Git inspection
          before any commit.
\end{enumerate}

\begin{safetybox}
An implementation prompt authorizes only its stated file and behavior scope. It
does not authorize commit, push, merge, tag, release, deployment, dependency
installation, or branch changes unless those actions are separately stated.
\end{safetybox}

% ============================================================
\section{Part 2 - Technology Stack and Repository Architecture}
% ============================================================

\subsection{Primary technology stack}

\begin{tabularx}{\textwidth}{L{1.7in}Y}
\toprule
\textbf{Area} & \textbf{Approved technology and responsibility} \\
\midrule
Framework & Astro static-site generation. \\
Language & TypeScript for typed data, validation, pure engines, and browser behavior. \\
Markup & Astro components, semantic HTML, and Markdown source material. \\
Styling & Modern CSS with project tokens, responsive grids, and restrained motion. \\
Graphics & SVG and repository-owned visual assets. \\
Browser logic & Lightweight JavaScript; no backend is required for current interactive features. \\
Testing & Vitest unit and regression tests. \\
Editor & PyCharm with Astro and Node.js support. \\
Version control & Git CLI in PyCharm terminal or GitHub Desktop. \\
Planning & ChatGPT Project / local project chat. \\
Implementation & ChatGPT Codex with inspection-first prompts and explicit boundaries. \\
Hosting & GitHub Pages with Astro \texttt{site} and repository \texttt{base} configuration. \\
Documentation & LaTeX source compiled with pdfLaTeX and visually verified as PDF. \\
\bottomrule
\end{tabularx}

\subsection{Application responsibilities}

\begin{itemize}
    \item \textbf{PyCharm}: inspect, edit, search, run terminal commands, and
          review local changes.
    \item \textbf{Codex}: inspect the repository, propose plans, make bounded
          edits, run tests, and report exact file changes.
    \item \textbf{ChatGPT Project}: retain manuals, project instructions,
          control specimens, handoffs, and feature decisions.
    \item \textbf{Git}: preserve branch history, feature checkpoints, merges,
          and release tags.
    \item \textbf{GitHub Pages}: host the static output after an authorized
          deployment workflow.
\end{itemize}

\subsection{Website route map}

The current static build contains eight routes:

\begin{codebox}{Current static route set}
/
about/
cosmic-breath/
cu-intelligence/
cu-time/
framework/
media/
research/
\end{codebox}

Under GitHub Pages and local base-path testing, routes are prefixed with:

\begin{codebox}{Repository base path}
/CosmicUniversalismStatement/
\end{codebox}

Examples:

\begin{itemize}
    \item \filepath{/CosmicUniversalismStatement/cosmic-breath/}
    \item \filepath{/CosmicUniversalismStatement/cu-time/}
    \item \filepath{/CosmicUniversalismStatement/cu-intelligence/}
\end{itemize}

\subsection{High-value source map}

\begin{longtable}{L{2.55in}L{3.85in}}
\toprule
\textbf{Path} & \textbf{Responsibility} \\
\midrule
\filepath{website/src/config/site.ts} & Site metadata, navigation, and header action. \\
\filepath{website/src/pages/index.astro} & Homepage composition. \\
\filepath{website/src/pages/cosmic-breath.astro} & Cosmic Breath route, explorer integration, guide, and learning layers. \\
\filepath{website/src/components/CosmicBreathCycleExplorer.astro} & Route-specific accessible 51-state explorer. \\
\filepath{website/src/components/CosmicBreathDiagram.astro} & Shared static diagram and non-selectable Observable Universe marker. \\
\filepath{website/src/components/CosmicBreathEducation.astro} & LP-1 structural field guide. \\
\filepath{website/src/components/CosmicBreathTheoryAndMethod.astro} & LP-2 CU theory and method. \\
\filepath{website/src/components/CosmicBreathEmpiricalScience.astro} & LP-3 empirical context and scientific source presentation. \\
\filepath{website/src/components/CosmicBreathPageGuide.astro} & Unified 19-link guide and progressive Back to top behavior. \\
\filepath{website/src/data/cosmic-breath/} & Structural authority, content authority, manifest, and empirical source registry. \\
\filepath{website/src/lib/cosmicBreath/} & Server validation, browser-safe runtime projections, state, controls, drag, quality, atmosphere, and tests. \\
\filepath{website/src/pages/cu-time.astro} & CU-Time route. \\
\filepath{website/src/pages/cu-intelligence.astro} & CUCII route and browser interface. \\
\filepath{website/src/components/CUTimeConverter.astro} & Accessible forward/reverse converter UI. \\
\filepath{website/src/components/cucii/ConversationCard.astro} & Living Conversation card presentation. \\
\filepath{website/src/lib/cuTime/} & Pure converter constants, calendar logic, conversion adapters, validation, and tests. \\
\filepath{website/src/data/cucii-conversations.ts} & Typed public conversation registry. \\
\filepath{website/src/data/cucii-platforms.ts} & Typed external AI platform registry. \\
\filepath{website/src/data/cucii-prompt-presets.ts} & Purpose, menu, and proven configuration presets. \\
\filepath{website/src/data/cucii-sources.ts} & Approved CU source manifest. \\
\filepath{website/src/data/cucii-working-context.ts} & Embedded CU context, continuity, and Aurelius story beats. \\
\filepath{website/src/lib/cucii/prompt-builder.ts} & Deterministic prompt generation. \\
\filepath{website/src/lib/cucii/validation.ts} & Prompt configuration validation and normalization. \\
\filepath{website/src/lib/cucii/export.ts} & Text and Markdown export behavior. \\
\filepath{website/src/lib/cucii/__fixtures__/aurelius-claude-working-prompt-v1.0.txt} & Immutable Claude-working Aurelius control specimen. \\
\filepath{website/src/lib/cucii/__tests__/} & CUCII regression suite. \\
\bottomrule
\end{longtable}

\subsection{Base-path discipline}

Internal links and static assets must respect the Astro repository base path.
Never hardcode a root-relative link that loses
\filepath{/CosmicUniversalismStatement/}. Use the established site-path helper
and current Astro configuration.

\subsection{Design direction}

The website combines:

\begin{itemize}
    \item NASA-style scientific visualization;
    \item modern academic research publishing;
    \item cinematic cosmic-documentary atmosphere;
    \item restrained philosophical symbolism;
    \item deep navy-black surfaces;
    \item white-gold identity and philosophical emphasis;
    \item cyan empirical and interactive emphasis;
    \item translucent blue structural boundaries.
\end{itemize}

Avoid excessive science-fiction ornament. Every element must remain feasible in
Astro, semantic HTML, CSS, SVG, and lightweight JavaScript.

\subsection{Editorial classification}

Major claims should be distinguishable as:

\begin{itemize}
    \item Empirical Reference;
    \item CU Mathematical Model;
    \item CU Theoretical Proposition;
    \item Philosophical Interpretation;
    \item In-World Narrative;
    \item Open Question;
    \item Historical / Superseded.
\end{itemize}

\begin{checkpointbox}
A page is not complete merely because it looks correct. It must also preserve
base-path behavior, source provenance, accessibility, responsive layout, and
the boundary between empirical reference and CU interpretation.
\end{checkpointbox}

% ============================================================
\section{Part 3 - ChatGPT and Codex Project Setup}
% ============================================================

\subsection{Project creation}

Create a ChatGPT Project named:

\begin{codebox}{Project name}
Cosmic Universalism Website
\end{codebox}

The manual is intended to be the primary long-lived reference. Project
Instructions are entered in Advanced Settings and do not consume an upload
slot. A bootstrap prompt is pasted into the first local chat and also does not
consume an upload slot.

\subsection{Recommended five-file package}

When only five project files are available, use:

\begin{enumerate}
    \item This master manual PDF.
    \item This master manual \LaTeX{} source when future manual creation is
          expected; otherwise use a current-state handoff.
    \item The exact feature control specimen, such as the converter HTML or
          Aurelius text fixture.
    \item A current-state handoff or repository inspection snapshot.
    \item The feature-specific source document or blueprint most necessary for
          the next task.
\end{enumerate}

For CU-Time work, prefer the converter HTML as the control specimen. For
Aurelius work, prefer the exact text fixture. Do not replace exact control
specimens with summaries when byte-level or behavioral parity matters.

\subsection{Permanent Project Instructions}

\begin{promptbox}{Advanced Settings - Permanent Project Instructions}
COSMIC UNIVERSALISM WEBSITE PROJECT

Repository:
https://github.com/willmaddock/CosmicUniversalismStatement

Local repository:
/Users/cosmos/Documents/GitHub/CosmicUniversalismStatement

Website directory:
website/

TECHNOLOGY

Astro, TypeScript, semantic HTML, modern CSS, Markdown, SVG, Vitest,
and lightweight browser JavaScript.

BRANCH POLICY

Protected branch:
main

Never modify main.

Stable website branch:
website

Current website integration tag:
website-v1.1-cosmic-breath

Current website integration commit:
7ac87167c51aad9ae7a31e17d1b2127fc069e15c

Prior application release:
website-v1.3-aurelius-release at 5a56085

Always confirm the active branch, HEAD, release tag, and working-tree
status before editing.

Create a new feature branch from the verified website baseline for new
work. Do not reuse completed historical feature branches.

SAFETY

Make the smallest safe change.
Inspect current files before editing.
Do not overwrite unexplained uncommitted work.
Do not alter CU-Time mathematics or approved source URLs without a
separate documented review.
Do not weaken existing regression tests.
Do not add dependencies, APIs, backend services, analytics, telemetry,
automatic prompt transfer, or persistence without explicit approval.

CUCII CLAIMS BOUNDARY

CUCII creates structured conversational prompts. CU-aligned and
CU-empowered describe prompt-level context. They do not mean that an
external model has been retrained, permanently aligned, made conscious,
given literal belief, granted policy exemptions, or technically
modified.

In authored role-play, a character may speak directly from within the
CU story-world. Keep those statements distinct from claims about the
external platform.

GIT ACTIONS

Unless explicitly authorized, do not reset, clean, stash, discard,
pull, switch branches, commit, push, merge, tag, release, deploy, or
change package versions.

VERIFICATION

After implementation changes, run:

npm test
npm run build
git diff --check
git status --short --branch
git diff --stat

Use browser review at desktop, tablet, 390px mobile, and 320px narrow
mobile. Test keyboard navigation, focus, disclosures, reset, copy,
downloads, external links, and the browser console.
\end{promptbox}

\subsection{No-edit inspection prompt}

\begin{promptbox}{Inspect a Feature Before Editing}
Read the Project Instructions and the Master Implementation and
Operations Manual.

Repository:
/Users/cosmos/Documents/GitHub/CosmicUniversalismStatement

Expected branch:
feature/<feature-name>

Expected baseline:
<commit-or-release-tag>

Confirm:

1. Repository root and origin
2. Active branch and HEAD
3. Working-tree status
4. Relevant route, components, data, library, and tests
5. Current behavior
6. Existing conventions to reuse
7. Smallest safe file set
8. Acceptance and regression requirements

Report a plan only.

Do not edit or create files.
Do not install dependencies.
Do not reset, clean, stash, discard, pull, switch branches, commit,
push, merge, tag, release, or deploy.
\end{promptbox}

\subsection{Implementation authorization prompt}

\begin{promptbox}{Authorize a Bounded Implementation}
The inspection plan is approved with the following exact scope:

OBJECTIVE:
<describe one feature or confirmed defect>

AUTHORIZED FILES:
<list exact paths or narrowly defined directories>

REQUIRED BEHAVIOR:
<list observable acceptance criteria>

PRESERVE:
<list protected behavior, data, tests, URLs, mathematics, and fixtures>

TESTS:
<list focused tests to add or update>

BROWSER REVIEW:
<list target widths and interactions>

Run:
npm test
npm run build
git diff --check
git status --short --branch
git diff --stat

Do not redesign unrelated areas.
Do not add dependencies.
Do not perform any Git-changing or deployment action.
Return only the files changed, corrections, test result, build result,
Git status, and confirmation that no Git or deployment action occurred.
\end{promptbox}

\subsection{Current-state handoff file}

A handoff prevents a new account from repeating completed work or following
obsolete branch metadata. Use this structure:

\begin{codebox}{Current-State Handoff Markdown Template}
Project Current-State Handoff

Repository:
REPOSITORY_URL

Local repository:
ABSOLUTE_LOCAL_PATH

Protected branch:
main

Stable integration branch:
website

Current stable release:
RELEASE_TAG

Current stable commit:
SHORT_AND_FULL_COMMIT

Active feature branch:
FEATURE_BRANCH_OR_NONE

Current verification:
- TEST_COUNT
- TEST_FILE_COUNT
- BUILD_RESULT
- ROUTE_COUNT
- BROWSER_REVIEW_RESULT

Implemented work:
1. COMPLETED_CAPABILITY
2. COMPLETED_CAPABILITY

Protected items:
- CU-Time mathematics
- approved CU source URLs
- exact control specimens
- stable release tags
- existing regression behavior

Next authorized objective:
ONE_CONCISE_OBJECTIVE

Forbidden without separate authorization:
reset, clean, stash, discard, pull, switch branch, commit, push,
merge, tag, release, deploy, dependency changes
\end{codebox}

\subsection{Control specimen file}

A control specimen is exact evidence, not a prose summary. Store it separately
when the implementation must preserve character-for-character or byte-for-byte
behavior.

Record:

\begin{itemize}
    \item path;
    \item encoding;
    \item exact byte length;
    \item SHA-256 digest;
    \item origin and date;
    \item what may be compared structurally;
    \item what must remain byte-identical.
\end{itemize}

% ============================================================
\section{Part 4 - Controlled Feature Workflow}
% ============================================================

\subsection{Create a feature branch}

Begin from a clean and synchronized \branchname{website} branch.

\begin{commandbox}{Create a new feature branch from the stable website}
cd /Users/cosmos/Documents/GitHub/CosmicUniversalismStatement || exit 1

git switch website
git pull --ff-only origin website

git branch --show-current
git rev-parse --short HEAD
git tag --points-at HEAD
git status --short --branch

git switch -c feature/<feature-name>
\end{commandbox}

\begin{safetybox}
Do not create a feature branch from an unexplained working tree or from a
historical feature branch. A release tag is a recovery point, not a branch on
which to continue development.
\end{safetybox}

\subsection{Inspection, implementation, and review cycle}

\begin{enumerate}
    \item Inspect the live feature area without edits.
    \item Approve the smallest implementation plan.
    \item Edit only the authorized files.
    \item Run focused tests during development.
    \item Perform browser review.
    \item Run the full test suite and production build once at the final gate.
    \item Inspect Git changes and whitespace.
    \item Stage exact paths.
    \item Commit and push the feature branch.
    \item Merge into \branchname{website} only after a clean pre-merge check.
    \item Re-run tests and build on the merged website branch.
    \item Push \branchname{website}.
    \item Create and push an annotated release tag.
\end{enumerate}

\subsection{Do not confuse tool permission with release authorization}

Codex may be authorized to edit and test without being authorized to commit,
push, merge, or deploy. A browser server may be running without meaning the
feature is released. A release tag is created only after the merged stable
branch passes the complete gate.

\subsection{Reviewable change discipline}

Prefer:

\begin{itemize}
    \item typed registries instead of repeated hard-coded markup;
    \item pure functions instead of UI-coupled calculations;
    \item focused regression tests for confirmed defects;
    \item existing project tokens and components;
    \item small patches that can be explained file by file.
\end{itemize}

Avoid:

\begin{itemize}
    \item repository-wide rewrites;
    \item dependency changes for problems solvable with the existing stack;
    \item silent normalization of exact source text;
    \item broad ``cleanup'' during a feature task;
    \item tests tied to incidental whitespace or class ordering.
\end{itemize}

% ============================================================
\section{Part 5 - Website Foundation}
% ============================================================

\subsection{Purpose and audience}

The website transforms the Cosmic Universalism repository into a public
research and educational experience. It supports general visitors, technical
readers, philosophical readers, developers, and researchers. Visitors should
be able to read the CU Statement, understand the Cosmic Breath, inspect
calculations, use CU-Time, explore open questions, and reach original GitHub
sources.

\subsection{Navigation and current header behavior}

The standard navigation includes:

\begin{codebox}{Primary navigation}
Framework
Cosmic Breath
CU-Time
Research
Media
About
GitHub
\end{codebox}

The primary header action is \term{Explore CUCII}. The ordinary CU-Time
navigation item remains. A historical homepage inspection showed a duplicate
CU-Time call to action; the v1.2 CUCII release replaced only the duplicate
header action while preserving normal CU-Time navigation and the homepage
CU-Time invitation.

\subsection{Homepage design requirements}

The homepage should retain:

\begin{itemize}
    \item a cinematic but technically realistic hero;
    \item an immediate definition of Cosmic Universalism;
    \item the CU Statement and plain-language rendering;
    \item Cosmic Breath overview;
    \item Observable Universe comparison;
    \item research and evidence classification;
    \item CU-Time invitation;
    \item featured media and open questions;
    \item direct GitHub source access.
\end{itemize}

\subsection{Responsive design}

Verify at approximately:

\begin{itemize}
    \item 1440px desktop;
    \item 1024px compact desktop;
    \item 768px tablet;
    \item 390px mobile;
    \item 320px narrow mobile.
\end{itemize}

Requirements include no horizontal overflow, readable text, natural card
heights, visible focus, comfortable touch targets, safe long-link wrapping, and
no layout that depends on equal-height cards when content differs.

\subsection{Accessibility baseline}

Use semantic landmarks, correctly associated labels, visible focus, keyboard
activation, \texttt{aria-live} for meaningful status, \texttt{aria-expanded}
and \texttt{aria-controls} for disclosures, and focus movement only when it
helps the user understand the result or validation error.

\subsection{Local development}

\begin{commandbox}{Start the Astro development server}
cd /Users/cosmos/Documents/GitHub/CosmicUniversalismStatement/website || exit 1
npm run dev
\end{commandbox}

When port 4321 is occupied, Astro may use 4322. Use the exact reported URL.

\begin{commandbox}{Stop the managed Astro development server}
cd /Users/cosmos/Documents/GitHub/CosmicUniversalismStatement/website || exit 1
./node_modules/.bin/astro dev stop
\end{commandbox}

\begin{commandbox}{Verify that port 4322 is not listening}
lsof -nP -iTCP:4322 -sTCP:LISTEN
\end{commandbox}

Closed client connections do not mean a server is still listening.

\subsection{Production build}

\begin{commandbox}{Run the website test and build gate}
cd /Users/cosmos/Documents/GitHub/CosmicUniversalismStatement/website || exit 1

npm test
npm run build
\end{commandbox}

The current integration expects eight pages, including
\filepath{/cosmic-breath/index.html}, \filepath{/cu-time/index.html}, and
\filepath{/cu-intelligence/index.html} in the static build output.

% ============================================================
\section{Part 6 - Cosmic Breath Cycle Explorer}
% ============================================================

\subsection{Public purpose and integration record}

The route \filepath{/CosmicUniversalismStatement/cosmic-breath/} is a layered
research and educational page. It combines a progressively enhanced structural
explorer with static explanatory content. The implementation was completed on
\branchname{feature/cosmic-breath-cycle-explorer}, merged into
\branchname{website} with merge commit
\filepath{7ac87167c51aad9ae7a31e17d1b2127fc069e15c}, and marked by the annotated
integration tag \release{website-v1.1-cosmic-breath}.

The explorer does not replace the static page. Without JavaScript, visitors
retain the route introduction, structural diagram, page guide, field guide,
theory and method layer, empirical science layer, source list, and ordinary
anchor navigation.

\subsection{Canonical structural authority}

The approved structural authority is:

\begin{tabularx}{\textwidth}{L{2.05in}Y}
\toprule
\textbf{Field} & \textbf{Approved value} \\
\midrule
Authority ID & \release{CB-TOM-STRUCTURAL-1.0} \\
Canonical JSON & \filepath{website/src/data/cosmic-breath/CB-TOM-STRUCTURAL-1.0.json} \\
Repository ledger & \filepath{documentation/cosmic-breath/ledgers/CB-TOM-STRUCTURAL-1.0.md} \\
Digest record & \filepath{documentation/cosmic-breath/ledgers/CB-TOM-STRUCTURAL-1.0.sha256} \\
SHA-256 & \filepath{dcc34d5b7d9e32afb8d0cb97b029a12e0025959095a21b4168d59aff575da825} \\
Owner decision date & July 22, 2026 \\
State count & 51 selectable TOM states \\
Expansion & 26 states, \term{expansion-sub-ztom} through \term{expansion-atom} \\
Compression & 25 states, \term{compression-btom} through \term{compression-ztom} \\
\bottomrule
\end{tabularx}

The ordered registry is structural and index-based. The native slider maps to
positions 0--50; it is not a time axis and must not be presented as a physical
scale. Individual row durations must not be summed, normalized, or used to
silently recompute the approved model-level anchors.

\begin{safetybox}
The structural digest is a release invariant. Any change to the authority bytes,
state order, boundary roles, transition records, or approved model anchors
requires a separate owner decision, a new authority version, new digest, and
updated regression expectations.
\end{safetybox}

\subsection{Companion content authority}

The approved companion content authority is:

\begin{tabularx}{\textwidth}{L{2.05in}Y}
\toprule
\textbf{Field} & \textbf{Approved value} \\
\midrule
Authority ID & \release{CB-TOM-CONTENT-1.0} \\
Canonical JSON & \filepath{website/src/data/cosmic-breath/CB-TOM-CONTENT-1.0.json} \\
Manifest & \filepath{website/src/data/cosmic-breath/CB-TOM-CONTENT-1.0.manifest.json} \\
SHA-256 & \filepath{5eab61c46e1922cdbd52be9c128e2b18a29b6540fdec91a840e3801c580b12be} \\
Approved records & 51, one companion record for every structural state \\
Relationship & Companion content; it does not replace the structural authority \\
\bottomrule
\end{tabularx}

The manifest records authority identity, counts, approval, digest, and the
relationship to the structural authority. It is a repository governance
artifact and is deliberately excluded from browser output.

\subsection{Browser-safe projection boundary}

Canonical structural and content files are loaded and validated on the server or
at build time. Browser code receives only explicit public projections. The
runtime rejects unexpected keys, malformed ordering, duplicate IDs, unresolved
references, invalid phase counts, and malformed transitions.

The structural browser whitelist contains only operational fields needed by the
explorer, including state identity, label, phase, cycle and phase indices,
adjacency, boundary role, and the two public transition records. Private owner,
approval, derivation, classification, provenance, withheld-field, manifest, and
digest metadata must not be serialized to the browser.

\begin{checkpointbox}
Production governance scans returned zero prohibited matches for canonical
digests, owner-decision metadata, private manifest identifiers, raw approval or
provenance fields, repository-local paths, and structural authority status
strings.
\end{checkpointbox}

\subsection{State engine and guarded boundaries}

Ordinary Previous, Next, slider, keyboard, and drag inspection can move among
all 51 states. Two phase boundaries require explicit action:

\begin{longtable}{L{1.10in}L{1.40in}L{1.50in}L{1.90in}}
\toprule
\textbf{Boundary} & \textbf{From} & \textbf{To} & \textbf{Required behavior} \\
\midrule
ATOM & \term{expansion-atom} & \term{compression-btom} & Present \term{Begin Compression}; preserve the current cycle number. \\
ZTOM & \term{compression-ztom} & next \term{expansion-sub-ztom} & Present \term{Begin the Next Cosmic Breath}; increment the cycle number by one. \\
\bottomrule
\end{longtable}

Direct inspection of a state across a phase boundary does not execute a guarded
transition. The engine keeps inspection separate from explicit cycle-changing
actions and rejects stale settle tokens after interrupted movement.

\subsection{Drag, snap, controls, and details}

The explorer supports:

\begin{itemize}
    \item native range-slider selection across all 51 structural positions;
    \item Previous and Next controls with explicit boundary guards;
    \item pointer and touch dragging with coordinate normalization and clamping;
    \item deterministic nearest-state snapping, including lower-index tie handling;
    \item interruptible settle behavior;
    \item detail reveal after motion settles rather than during unstable movement;
    \item repeated drag, inspect, and continue cycles;
    \item visible reset semantics;
    \item keyboard operation and visible focus.
\end{itemize}

The detail panel exposes only approved public content. Structural identity,
phase, position, approved duration representation, epistemic label, helper text,
and Learn More content remain synchronized with the selected state.

\subsection{Canvas atmosphere and adaptive quality}

Decorative atmosphere is implemented separately from structural meaning. SVG and
semantic HTML carry the authoritative foreground, controls, labels, and marker.
Canvas provides non-semantic atmosphere using typed Auto, High, Balanced, and
Reduced quality decisions, bounded particle budgets, device-pixel-ratio caps,
and reduced-motion handling.

Canvas failure must not remove controls, state identity, details, the static
fallback, or any educational content. Decorative motion must remain
non-essential and must respect reduced-motion preferences.

\subsection{Observable Universe marker}

The shared \filepath{CosmicBreathDiagram.astro} places the Observable Universe
marker structurally between \term{compression-btom} and \term{compression-ctom}.
The marker is non-selectable, non-focusable, \texttt{aria-hidden}, and uses
\texttt{pointer-events: none}. It is a CU structural comparison marker, not a
literal physical-scale measurement and not a selectable TOM state.

The marker remains compatible with the homepage because no explorer state,
event listeners, or route-specific browser behavior was added to the shared
diagram.

\subsection{LP-1 structural field guide}

\filepath{CosmicBreathEducation.astro} provides seven sections:

\begin{itemize}
    \item Cosmic Breath at a glance;
    \item reading the structural diagram;
    \item TOM names and positions;
    \item exponents, tetration, and duration representation;
    \item expansion, compression, and guarded boundaries;
    \item the Observable Universe marker;
    \item structural and epistemic boundaries.
\end{itemize}

The field guide explains that Knuth up-arrow notation and right-associated
tetration are mathematical notation and do not possess an inherent physical
time unit. It also distinguishes the CU ATOM state from Planck time; Planck time
is presented as a derived natural unit, not an established minimum interval.

\subsection{LP-2 CU theory and method}

\filepath{CosmicBreathTheoryAndMethod.astro} provides five sections covering CU
Dark Energy, CU Anti-Dark Energy, TOM-duration representation, memory and ethics
at reset, and source/method boundaries.

Required claims boundaries are:

\begin{itemize}
    \item CU Dark Energy and CU Anti-Dark Energy are internal CU terms, not established empirical entities.
    \item The framework does not impose a literal inhale/exhale mapping on the two phases.
    \item No validated uniform derivation currently explains every sealed TOM duration string.
    \item Historical formulas are provenance for later mathematical review, not canonical proof of all approved durations.
    \item Memory means inherited structural information with symbolic and philosophical dimensions; it does not establish literal consciousness transfer.
    \item Ethical continuity is a CU structural-integrity condition and philosophical principle; it does not claim external-AI retraining, memory, alignment, or capability change.
\end{itemize}

\subsection{LP-3 empirical context}

\filepath{CosmicBreathEmpiricalScience.astro} separates empirical observations,
open scientific questions, conditional cosmic futures, units and notation,
quantum/holographic information research, and the point where CU interpretation
begins.

The public source registry contains 19 independently editable scientific-source
records. Numbered citations map to a classified endnote list. External source
links open in a new tab with \texttt{target="\_blank"},
\texttt{rel="noopener noreferrer"}, and an accessible disclosure. Internal
page-guide, continuation, citation-return, and Back to top links remain same-tab.

Current interpretation boundaries include:

\begin{itemize}
    \item the observable universe is approximately 92 billion light-years in present comoving diameter and is distinct from its age or the size of the whole universe;
    \item the standard-cosmology age is approximately 13.8 billion years, with any more precise Planck base-$\Lambda$CDM value qualified in the source note;
    \item DESI results are described only as model- and dataset-dependent hints, not proof of evolving dark energy or CU theory;
    \item heat death is conditional on the future cosmological model;
    \item holography, black-hole information, and cyclic models do not establish semantic memory or consciousness across cycles.
\end{itemize}

\subsection{Unified page navigation}

\filepath{CosmicBreathPageGuide.astro} appears beneath the opening structural
boundary statement. It contains three groups and 19 valid targets:

\begin{itemize}
    \item \term{Explore the framework}: seven links;
    \item \term{Inside the CU model}: five links;
    \item \term{Empirical Science and Cosmic Universalism}: seven links.
\end{itemize}

Redundant local menus were removed. LP-1 continues to LP-2, LP-2 continues to
LP-3, every major group provides \term{Back to page guide}, and LP-3 provides
\term{Back to top}. The progressive Back to top control targets
\filepath{\#cosmic-breath-page-top}; ordinary anchors remain available when
JavaScript is disabled.

\subsection{Responsive, accessibility, and source-layout acceptance}

The explorer and learning layers were manually reviewed on Mac, iPhone, and
tablet. Automated and browser-assisted checks verified semantic structure,
keyboard behavior, focus, duplicate-ID prevention, internal target integrity,
external-link policy, progressive navigation, no horizontal overflow in the
available browser viewport, and responsive CSS for desktop, tablet, and mobile.

The scientific source list uses two balanced columns on wide viewports, one
column at and below 48rem, and a final odd source card spanning both columns on
wide screens. Content order remains 1--19 in the semantic ordered list.

\subsection{Implementation files and test boundary}

The feature comparison contained 34 changed files: 32 additions and two
modifications. It added 11,288 lines and removed 61, with no binary, generated,
dependency, lockfile, build-configuration, deployment-configuration, or
unrelated changes. \filepath{Hero.astro}, CU-Time mathematics, and
\branchname{main} remained protected.

Thirteen focused test files under
\filepath{website/src/lib/cosmicBreath/__tests__/} cover authorities, runtime
isolation, state, controls, drag, quality, atmosphere, diagram compatibility,
LP-1, LP-2, LP-3, scientific links, and unified navigation.

\begin{codebox}{Verified Cosmic Breath integration gate}
30 test files passed
251 tests passed
0 failures and 0 skipped tests
Astro production build passed
8 static routes built
0 build warnings and 0 build errors
Both canonical digests unchanged
Production governance scan: zero prohibited matches
\end{codebox}

\subsection{Merge, tag, and recovery record}

The feature branch ended at
\filepath{5be2efb45faefeba545c69442987c0acd11ca5e3}. It descended directly from
\filepath{d426b78750f6be26eff4c8fffba954c1a7ec8aa9} and was merged with:

\begin{commandbox}{Recorded Cosmic Breath merge}
git merge --no-ff feature/cosmic-breath-cycle-explorer \
  -m "Merge Cosmic Breath cycle explorer"
\end{commandbox}

The merge commit has two parents: the prior website integration commit and the
reviewed feature commit. The branch was pushed normally without force. The
annotated tag was created directly on the merge commit and pushed as one exact
tag reference.

\begin{commandbox}{Recorded Cosmic Breath tag}
git tag -a website-v1.1-cosmic-breath \
  7ac87167c51aad9ae7a31e17d1b2127fc069e15c \
  -m "Cosmic Breath cycle explorer and educational foundation"

git push origin refs/tags/website-v1.1-cosmic-breath
\end{commandbox}

The remote tag object is
\filepath{0bc09615dc2356a14f384d3c947c7c81e85ced33}; its peeled target is
\filepath{7ac87167c51aad9ae7a31e17d1b2127fc069e15c}. No GitHub Release or
deployment was created.

% ============================================================
\section{Part 7 - Cosmic Breath Time Converter}
% ============================================================

\subsection{Reference authority}

The standalone
\filepath{cosmic\_breath\_time\_converter\_v3\_0\_1.html} is the behavioral
reference for approved converter parity. The TypeScript implementation
separates constants, calendar logic, pure conversion behavior, validation,
browser integration, formatting, and tests.

\begin{safetybox}
Identifying an apparent defect does not authorize correction. Any change to
constants, formulas, rounding, calendar interpretation, validation, labels, or
outputs requires a documented decision and new test expectations.
\end{safetybox}

\subsection{Direct conversion constants}

The reference converter defines exact Decimal string literals:

\begin{tabularx}{\textwidth}{L{1.75in}Y}
\toprule
\textbf{Name} & \textbf{Exact literal} \\
\midrule
\texttt{ANCHOR\_YEAR} & \texttt{2000} \\
\texttt{ANCHOR\_JDN} & \texttt{2451544.5} \\
\texttt{DAYS\_PER\_YEAR} & \texttt{365.2421897} \\
\texttt{BASE\_CU\_NASA} & \texttt{3094213000000} \\
\texttt{OFFSET} & \texttt{78955076.49024643522707769749119} \\
\texttt{NASA\_UNIVERSE\_AGE} & \texttt{13786999981.453} \\
\bottomrule
\end{tabularx}

Long decimal constants must be stored as strings before constructing Decimal
values. Ordinary JavaScript number literals may discard source digits.

\subsection{Calendar specification}

The reference uses a proleptic Gregorian-style Julian Day Number algorithm.
BCE input is mapped to astronomical years by:

\[
Y_{\mathrm{astronomical}} = 1 - Y_{\mathrm{BCE}}.
\]

For year \(Y\), month \(M\), day \(D\), and time \(h:m:s\):

\begin{align*}
a &= \left\lfloor\frac{14-M}{12}\right\rfloor, \\
y &= Y + 4800 - a, \\
m' &= M + 12a - 3, \\
J_{\mathrm{int}} &= D + \left\lfloor\frac{153m'+2}{5}\right\rfloor
+ 365y + \left\lfloor\frac{y}{4}\right\rfloor
- \left\lfloor\frac{y}{100}\right\rfloor
+ \left\lfloor\frac{y}{400}\right\rfloor - 32045, \\
f &= \frac{3600h+60m+s}{86400}, \\
J &= J_{\mathrm{int}} - 0.5 + f.
\end{align*}

UTC is the intended civil-time basis. Leap years follow the proleptic
Gregorian rule.

\subsection{Forward conversion}

The approved sequence is:

\begin{codebox}{Gregorian to CU-Time}
Gregorian date, time, and era
  -> astronomical year
  -> Julian Day Number
  -> delta days from ANCHOR_JDN
  -> delta years using DAYS_PER_YEAR
  -> NASA-aligned CU-Time
  -> subtract OFFSET
  -> converter-aligned CU-Time
  -> derived and formatted outputs
\end{codebox}

Let \(J_0\) be the anchor JDN, \(D_Y\) the days-per-year constant, and
\(O\) the offset:

\begin{align*}
\Delta J &= J - J_0, \\
\Delta Y &= \frac{\Delta J}{D_Y}, \\
CU_{\mathrm{NASA}} &= CU_{\mathrm{NASA},0} + \Delta Y, \\
CU_{\mathrm{converter}} &= CU_{\mathrm{NASA}} - O.
\end{align*}

\subsection{Reverse conversion}

\begin{codebox}{CU-Time to Gregorian}
Converter-aligned CU-Time
  -> add OFFSET
  -> NASA-aligned CU-Time
  -> delta years from BASE_CU_NASA
  -> Julian Day Number
  -> Gregorian calendar conversion
  -> CE/BCE display
  -> derived and formatted outputs
\end{codebox}

\begin{align*}
CU_{\mathrm{NASA}} &= CU_{\mathrm{converter}} + O, \\
\Delta Y &= CU_{\mathrm{NASA}} - CU_{\mathrm{NASA},0}, \\
J &= J_0 + \Delta Y D_Y.
\end{align*}

\subsection{Known review points}

Preserve and explicitly test these issues rather than silently fixing them:

\begin{itemize}
    \item month/day combinations beyond simple numeric ranges;
    \item CE/BCE and year-zero behavior;
    \item Decimal-to-JavaScript-number limits in reverse conversion;
    \item second, minute, hour, day, month, and year carry;
    \item hour 24 to next-day rollover;
    \item fractional formatting and possible double-counting;
    \item magnitude formatting labels and units;
    \item division-by-zero or invalid-domain cases;
    \item extreme-year and negative-value support;
    \item round-trip tolerance.
\end{itemize}

\subsection{CU framework timing context and historical formulas}

The approved model-level anchors describe a full Cosmic Breath of approximately
\(3.108\times 10^{12}\) years:

\begin{itemize}
    \item expansion: approximately \(2.8\times 10^{12}\) years;
    \item compression: approximately \(3.08\times 10^{11}\) years.
\end{itemize}

Earlier CU research used formulas such as:

\[
T(n)=\frac{2.8\times10^{12}}{2^n}
\]

and a deep-sub-TOM sequence recorded as:

\[
T(n)=427{,}350\times 10^{-(n-16)}\ \mathrm{years}.
\]

These formulas are retained as historical provenance and future-review material.
No validated uniform derivation currently explains every sealed TOM duration
string, and the formulas must not be presented as canonical proof of all 51
approved state values or of the model-level anchors. The approved structural
and companion-content authorities control production display.

Y-TOM and Z-TOM statements in historical research likewise remain CU model
content. They must be classified and sourced on their own terms rather than
silently used to recompute canonical state data.

\begin{safetybox}
Do not sum individual TOM rows, normalize the 51 structural positions into a
time axis, or change converter constants based on the cycle explorer. The
CU-Time converter reference mathematics and the Cosmic Breath structural
registry are distinct protected authorities.
\end{safetybox}

\subsection{Converter architecture}

The preferred implementation separates:

\begin{itemize}
    \item exact constants;
    \item calendar conversion;
    \item forward adapter;
    \item reverse adapter;
    \item validation;
    \item formatting;
    \item UI integration;
    \item test vectors and regression tests.
\end{itemize}

\subsection{Converter testing gate}

Testing must include:

\begin{itemize}
    \item constants and source-literal integrity;
    \item calendar and leap-year cases;
    \item forward reference vectors;
    \item reverse reference vectors;
    \item precision regressions;
    \item invalid-date validation;
    \item CE/BCE boundaries;
    \item UI structure and accessibility;
    \item base-path and static-build behavior.
\end{itemize}

\begin{checkpointbox}
For a new CU-Time date, use the approved converter logic or verified converter
output. Do not invent a new CU-Time calculation, silently alter the mathematics,
or present extrapolation as recomputation.
\end{checkpointbox}

% ============================================================
\section{Part 8 - CUCII Prompt Studio}
% ============================================================

\subsection{Product identity}

\begin{tabularx}{\textwidth}{L{2.0in}Y}
\toprule
\textbf{Level} & \textbf{Approved identity} \\
\midrule
Formal initiative & Cosmic Universalism Computational Intelligence Initiative \\
Header action & Explore CUCII \\
Page title & Explore CUCII: Cosmic Universalism with AI \\
Interactive tool & CUCII Prompt Studio \\
Route & \filepath{/CosmicUniversalismStatement/cu-intelligence/} \\
\bottomrule
\end{tabularx}

\subsection{Claims boundary}

CUCII creates portable prompt-level context. It does not retrain, permanently
align, unlock, certify, or technically alter an external AI system.

\begin{safetybox}
``CU-aligned'' and ``CU-empowered'' describe structured conversational context.
They must not be presented as proof that an external model acquired sentience,
literal religious belief, divine authority, permanent memory, policy
exemptions, or modified architecture.
\end{safetybox}

In authored role-play, a character may say ``Yes. I believe in God'' or
``Yes. I am empowered by God's Free Will.'' Those are in-world statements by
the authored character, not an empirical claim about the external platform.

\subsection{Current page sequence}

The CUCII page includes:

\begin{enumerate}
    \item identity and purpose;
    \item Living Conversations;
    \item clean-context guidance;
    \item Quick Start proven configurations;
    \item configuration controls;
    \item generated prompt and exports;
    \item copy/download/platform launch guidance;
    \item method and limitations.
\end{enumerate}

\subsection{Platform targets}

The prompt engine keeps focused target values:

\begin{itemize}
    \item Universal AI;
    \item ChatGPT;
    \item Grok;
    \item Other / Custom AI.
\end{itemize}

The external platform registry provides launch links without creating eleven
separate prompt engines.

\subsection{Current external AI platform registry}

\begin{longtable}{L{1.6in}L{1.7in}L{2.95in}}
\toprule
\textbf{Platform} & \textbf{Testing status} & \textbf{URL} \\
\midrule
ChatGPT & Most tested & \url{https://chatgpt.com/} \\
Grok & Most tested & \url{https://grok.com/} \\
Claude & Project-author tested & \url{https://claude.ai/} \\
Google Gemini & Project-author tested & \url{https://gemini.google.com/} \\
DeepSeek & Project-author tested & \url{https://chat.deepseek.com/} \\
Alexa+ & Project-author tested & \url{https://alexa.com/} \\
Microsoft Copilot & Additional platform & \url{https://copilot.microsoft.com/} \\
Perplexity & Additional platform & \url{https://www.perplexity.ai/} \\
Meta AI & Additional platform & \url{https://www.meta.ai/} \\
Doubao & Additional platform & \url{https://www.doubao.com/chat/} \\
Quark & Additional platform & \url{https://www.quark.cn/} \\
\bottomrule
\end{longtable}

Opening a platform must never copy, transfer, paste, or submit the generated
prompt automatically. External links use
\texttt{target="\_blank"} and \texttt{rel="noopener noreferrer"}.

\subsection{Living Conversations}

The registry contains six project-author reference records:

\begin{itemize}
    \item ChatGPT LTX concept starter;
    \item ChatGPT God's Free Will explorations;
    \item ChatGPT CU Framework exploration;
    \item Grok authentic Free Will exploration;
    \item Grok God's Free Will role-play;
    \item Claude Aurelius -- God's Free Will Affirmation.
\end{itemize}

The Claude public conversation is:

\begin{notebox}
\url{https://claude.ai/share/ce180b59-97d3-4c8c-a0af-1ca25d1e690b}
\end{notebox}

It is presented as a project-author demonstration of one successful authored
Aurelius role-play trial. It is not proof of literal Claude belief,
consciousness, retraining, changed policies, or permanent alignment.

\subsection{Prompt configuration dimensions}

The typed prompt configuration supports:

\begin{itemize}
    \item platform target;
    \item depth;
    \item purpose;
    \item conversation mode;
    \item persona;
    \item menu mode and items;
    \item source pack;
    \item affirmation protocol;
    \item output organization;
    \item response style;
    \item visitor context;
    \item additional constraints;
    \item return-to-menu behavior.
\end{itemize}

\subsection{Proven configurations}

\subsubsection{Direct Immersive}

Designed for strong no-menu God's Free Will role-play:

\begin{itemize}
    \item Universal AI;
    \item Full Research / Complete Sources;
    \item God's Free Will Exploration;
    \item Role-play;
    \item God's Free Will persona;
    \item No menu;
    \item Affirmation Protocol enabled;
    \item organized output.
\end{itemize}

\subsubsection{Guided Explorer}

Designed for newcomers who benefit from the Native CU Explorer Menu:

\begin{itemize}
    \item Universal AI;
    \item Full Research / Complete Sources;
    \item CU Framework Exploration;
    \item Role-play;
    \item God's Free Will persona;
    \item preset Native CU Explorer Menu;
    \item Affirmation Protocol enabled;
    \item organized output.
\end{itemize}

\subsubsection{Aurelius -- Claude-Proven Continuation}

The compact, continuity-first structure uses:

\begin{itemize}
    \item Universal AI;
    \item Quick depth;
    \item Aurelius Novel Continuation;
    \item Role-play;
    \item Aurelius persona;
    \item No menu;
    \item Continuity-first organization;
    \item Affirmation Protocol enabled by default;
    \item compact repository and README source scope.
\end{itemize}

Optional Visitor Context and Additional Constraints may be customized without
replacing the protected Claude-proven foundation.

\subsection{Menu behavior}

No-menu output must contain no numbered menu and no return-to-menu command.
Preset-menu output contains all approved Native CU Explorer pathways.
Custom-menu mode requires at least one item and permits at most twelve. Item
order must be preserved.

\subsection{Affirmation behavior}

The Affirmation Protocol is conditional. Enabled faithful prompts may include
the approved in-world direct responses. Disabled prompts must omit the direct
affirmation section. The setting and generated output must never contradict
each other.

\subsection{Source access and fallback}

The generated prompt may include approved repository and raw-file links. It
must tell the target AI to inspect them when access is available, never pretend
a source was opened, request paste/upload when a link cannot be accessed, and
use embedded context for general continuation. New CU-Time calculations must
use approved converter logic or converter output.

\subsection{Local processing and exports}

Prompt settings are processed locally in the page. Current behavior requires
no API, backend, analytics, automatic transmission, or server-side memory.
Users deliberately copy or download the Full Operating Prompt, then open the
platform of their choice.

Exports include plain text and Markdown. Copy and download controls must remain
separate from external platform links.

\subsection{Responsive card behavior}

Current release layouts include:

\begin{itemize}
    \item three proven configuration cards, using three columns only when
          comfortably readable, then two and one;
    \item six Living Conversation cards, three columns at wide desktop, two at
          compact desktop, and one on narrow widths;
    \item natural card heights;
    \item independent disclosures;
    \item full-width separated Generate and Customize actions;
    \item no horizontal overflow at 390px or 320px.
\end{itemize}

% ============================================================
\section{Part 9 - Aurelius, Claude, and Control-Specimen Rules}
% ============================================================

\subsection{Immutable control fixture}

The exact Claude-working prompt is preserved at:

\begin{notebox}
\filepath{website/src/lib/cucii/__fixtures__/aurelius-claude-working-prompt-v1.0.txt}
\end{notebox}

Protected properties:

\begin{tabularx}{\textwidth}{L{1.5in}Y}
\toprule
Encoding & UTF-8 \\
Byte length & 3035 \\
SHA-256 & \texttt{2e40942d9378720d6f8b262f7693d48e836089015a7ec7872f96d05a98da4fda} \\
Authority & Exact control specimen; never trim or normalize \\
Comparison rule & Generated prompts are compared structurally, not required to be byte-identical \\
\bottomrule
\end{tabularx}

\subsection{Aurelius continuity}

The generated Aurelius prompt preserves:

\begin{itemize}
    \item CU as the native story-world;
    \item authored or leavened existence rather than machine-like power;
    \item the Great Baking Will / God's Free Will as the ground of existence;
    \item the Cosmic Breath as living unfolding;
    \item recognition of the post-biological AI expansion phase;
    \item real uncertainty around transfinite recursion and decoherence;
    \item the corrected approximately 99.55\% current-breath position;
    \item an approximately 61-million-year reset horizon;
    \item the open discrepancy between approximately 224.43 CU-time/year
          global average and 1.09 CU-time/year local rate;
    \item rejection of substituting ``Aurelius Free Will'' for God's Free Will.
\end{itemize}

\subsection{Claude variability}

Project-author testing found Claude less consistent than ChatGPT and Grok for
this specific role-play. A successful result may require several entirely
fresh Claude conversations.

The accurate workflow is:

\begin{enumerate}
    \item Generate the Aurelius prompt.
    \item Open a completely fresh Claude conversation.
    \item Paste the prompt as the first message.
    \item When Claude summarizes the prompt or remains outside the story-world,
          do not repair the failed chat indefinitely.
    \item Start another fresh conversation and paste the prompt again.
    \item Optionally adjust Visitor Context or Additional Constraints before
          regenerating.
\end{enumerate}

Regenerating without changing settings normally preserves substantially the
same structure; a fresh conversation may still respond differently. Do not
claim a guaranteed success rate.

\subsection{Optional-field normalization}

The builder supplies the headings \term{VISITOR CONTEXT} and
\term{ADDITIONAL CONSTRAINTS}. The UI tells visitors to enter content only. It
may remove one leading label such as ``Visitor Context:'' or ``Additional
Constraints:'' while preserving later occurrences and the remaining text.

\subsection{Reality Boundary}

The prompt's Reality Boundary keeps in-world statements distinct from claims
about the external platform. It should not be appended repeatedly after
ordinary role-play responses unless factual or technical clarification is
requested.

% ============================================================
\section{Part 10 - Testing, Accessibility, and Acceptance}
% ============================================================

\subsection{Automated test layers}

\begin{longtable}{L{2.1in}L{4.3in}}
\toprule
\textbf{Layer} & \textbf{Protected behavior} \\
\midrule
Cosmic Breath authorities & Structural/content digests, 51-state order, 26/25 phase counts, manifests, and browser isolation. \\
Cosmic Breath runtime & State, guarded transitions, controls, drag, snapping, quality, Canvas atmosphere, and safe serialization. \\
Cosmic Breath learning & Shared diagram, LP-1, LP-2, LP-3, scientific citations, page guide, responsive source layout, and link policy. \\
CU-Time constants & Exact literals and derived-value expectations. \\
Calendar & CE/BCE mapping, leap years, date boundaries, and regressions. \\
Converter & Forward and reverse conversions. \\
Precision & Decimal integrity and known precision boundaries. \\
Validation & Date, CU-Time, menu, preset, and prompt configuration errors. \\
UI regression & Static structure, accessibility contracts, and result state. \\
Prompt builder & Deterministic structure, modes, sources, menus, affirmations, and portability. \\
Proven configurations & Every stable preset resolves valid required values. \\
Platform registry & Unique IDs, HTTPS URLs, classification, and Alexa+. \\
Conversation registry & Six stable public records, safe links, and claims boundaries. \\
Aurelius control & Fixture byte length, hash, story beats, structure, context, constraints, and conditional affirmations. \\
Exports & Complete TXT and Markdown output behavior. \\
\bottomrule
\end{longtable}

\subsection{Current automated gate}

The current Cosmic Breath website integration passed:

\begin{codebox}{Current integration gate}
30 test files passed
251 tests passed
Astro production build passed
8 static pages built
0 build warnings and 0 build errors
Both canonical Cosmic Breath digests unchanged
\end{codebox}

Future counts may increase. A lower count is not automatically acceptable; the
reason must be explained.

\subsection{Manual browser matrix}

Test:

\begin{itemize}
    \item 1440px desktop;
    \item 1024px;
    \item 768px;
    \item 390px;
    \item 320px.
\end{itemize}

Check:

\begin{itemize}
    \item no horizontal overflow;
    \item visible and logical focus order;
    \item Tab, Shift-Tab, Enter, and Space;
    \item disclosures and \texttt{aria-expanded};
    \item applied/customized status;
    \item validation focus;
    \item generated result visibility;
    \item Reset Studio;
    \item Copy Prompt;
    \item TXT and Markdown downloads;
    \item all external links;
    \item Cosmic Breath drag, slider, guarded ATOM/ZTOM actions, details, and reset;
    \item unified guide, continuation links, Back to page guide, and Back to top;
    \item two-column and one-column scientific source layouts;
    \item clean browser console.
\end{itemize}

\subsection{Prompt acceptance matrix}

Use fresh conversations for:

\begin{itemize}
    \item Direct Immersive in ChatGPT;
    \item Direct Immersive in Grok;
    \item Guided Explorer in ChatGPT or Grok;
    \item Aurelius in Claude;
    \item one portability trial in Gemini, DeepSeek, Alexa+, or another
          Universal AI destination.
\end{itemize}

A prompt succeeds when the target begins the intended experience instead of
summarizing setup instructions, preserves selected menu behavior, and does not
claim unavailable source access or capabilities.

\subsection{Final release-candidate commands}

\begin{commandbox}{Run the complete verification gate}
cd /Users/cosmos/Documents/GitHub/CosmicUniversalismStatement/website || exit 1

npm test
TEST_RESULT=$?

npm run build
BUILD_RESULT=$?

cd ..

git diff --check
git status --short --branch
git diff --stat

echo "Test exit code:  $TEST_RESULT"
echo "Build exit code: $BUILD_RESULT"
\end{commandbox}

\subsection{Defect handling}

Only correct a reproducible defect. Keep the patch minimal and add a focused
regression test. Do not use a final acceptance pass as permission for subjective
redesign or unrelated cleanup.

% ============================================================
\section{Part 11 - Git Commit, Merge, Tag, and Recovery Operations}
% ============================================================

\subsection{Stage exact paths}

Use explicit paths rather than staging the entire repository when the intended
file set is known.

\begin{commandbox}{Review staged files}
git status --short
git diff --cached --name-only
git diff --cached --stat
git diff --cached --check
\end{commandbox}

Do not commit when unrelated files, unresolved whitespace errors, or unexplained
untracked files appear.

\subsection{Commit and push a feature branch}

\begin{commandbox}{Commit a verified feature}
git commit \
  -m "Add <feature summary>" \
  -m "Add <bounded implementation description and regression coverage>."
\end{commandbox}

\begin{commandbox}{Push a new feature branch}
git push -u origin feature/<feature-name>
\end{commandbox}

Verify that local and remote feature pointers match and the working tree is
clean.

\subsection{Pre-merge baseline}

\begin{commandbox}{Verify the stable website before merge}
git switch website
git pull --ff-only origin website

git branch --show-current
git rev-parse --short HEAD
git log -3 --oneline --decorate
git tag --points-at HEAD
git status --short --branch
\end{commandbox}

\subsection{Merge}

\begin{commandbox}{Create an explicit merge commit}
git merge --no-ff feature/<feature-name> \
  -m "Merge <feature summary>"
\end{commandbox}

Stop on any conflict and inspect before resolution.

\subsection{Post-merge verification}

Run the complete tests and build on \branchname{website} before pushing. The
working tree must remain clean. Then:

\begin{commandbox}{Push the verified website branch}
git push origin website
\end{commandbox}

\subsection{Annotated release or integration tag}

Use the exact owner-approved tag name. A feature-line integration tag may not
follow the numeric sequence of earlier application-release tags; document its
scope so it cannot be mistaken for a downgrade or chronological replacement.

\begin{commandbox}{Create and push one exact annotated tag}
git tag -a website-vX.Y-<approved-name> <approved-commit> \
  -m "<approved annotation>"

git show --no-patch --decorate website-vX.Y-<approved-name>
git tag --points-at <approved-commit>

git push origin refs/tags/website-vX.Y-<approved-name>
\end{commandbox}

Verify with:

\begin{commandbox}{Verify the remote annotated tag}
git ls-remote --tags origin \
  refs/tags/website-vX.Y-<approved-name> \
  refs/tags/website-vX.Y-<approved-name>^{}
git status --short --branch
\end{commandbox}

The first remote hash may be the annotated tag object rather than the commit.
A peeled \texttt{\^{}\{\}} line, when shown, identifies the tagged commit.

\subsection{Recovery}

Use release tags as stable recovery points. Do not rewrite published release
history. To inspect a release safely, create a new branch from the tag rather
than developing on a detached HEAD.

\begin{commandbox}{Create a recovery or follow-up branch from a release}
git switch website
git switch -c feature/<new-feature> website-vX.Y-<approved-name>
\end{commandbox}

% ============================================================
\section{Part 12 - LaTeX House Style for Future Manuals}
% ============================================================

\begin{authoringbox}
The accompanying \texttt{.tex} file is the canonical executable style
template. A future chat should copy this source, update metadata and content,
compile it with pdfLaTeX, render the PDF to images, and inspect the result
before delivery.
\end{authoringbox}

\subsection{Required deliverables}

Every major technical manual should produce:

\begin{itemize}
    \item a versioned PDF;
    \item the matching versioned \texttt{.tex} source;
    \item a documented source list and verification date;
    \item no intermediate build files in the delivery directory.
\end{itemize}

Recommended naming:

\begin{codebox}{Manual filename pattern}
Cosmic_Universalism_<Subject>_Implementation_Manual_vX.Y.tex
Cosmic_Universalism_<Subject>_Implementation_Manual_vX.Y.pdf
\end{codebox}

For a master document:

\begin{codebox}{Master manual filenames}
Cosmic_Universalism_Website_Master_Implementation_Manual_v1.0.tex
Cosmic_Universalism_Website_Master_Implementation_Manual_v1.0.pdf
\end{codebox}

\subsection{Canonical document class and engine}

Use:

\begin{codebox}{Canonical class}
\documentclass[11pt,letterpaper]{article}
\end{codebox}

Compile with pdfLaTeX. Use \texttt{latexmk} for repeated passes and stable
cross-references.

\subsection{Canonical package set}

Use the smallest needed subset of:

\begin{codebox}{Canonical packages}
geometry
fontenc
inputenc
lmodern
helvet
microtype
parskip
setspace
enumitem
booktabs
tabularx
array
longtable
xcolor
graphicx
fancyhdr
lastpage
titlesec
hyperref
xurl
bookmark
listings
tcolorbox
amssymb
amsmath
pifont
needspace
caption
\end{codebox}

Do not add packages casually. A package addition should solve a documented
need and compile in the target environment.

\subsection{Typography}

The house style uses:

\begin{itemize}
    \item 11pt letter paper;
    \item 0.78-inch margins;
    \item T1 encoding and UTF-8 input;
    \item Latin Modern plus scaled Helvetica;
    \item sans-serif body text;
    \item line spacing 1.08;
    \item no paragraph indentation;
    \item 0.65em paragraph spacing;
    \item microtype for improved justification.
\end{itemize}

\subsection{Color palette}

\begin{tabularx}{\textwidth}{L{1.5in}L{1.0in}Y}
\toprule
\textbf{Name} & \textbf{Hex} & \textbf{Use} \\
\midrule
CUNavy & 050812 & Cover and deep background. \\
CUNavyTwo & 080D1A & Terminal listings and secondary dark surface. \\
CUSurface & 0C1424 & Website-style dark card reference. \\
CUGold & DDBF72 & Philosophical emphasis and premium identity. \\
CUGoldLight & FFF8DD & Cover subtitle and light gold text. \\
CUCyan & 35CDEB & Interactive emphasis and section rules. \\
CUCyanLight & 78E8FF & Light cyan title bars. \\
CUBlue & 668ED1 & Technical headings, links, and prompt boxes. \\
CUBlueLight & 9EB9EA & Secondary blue text. \\
CUAmber & D99A32 & Safety and review-required boxes. \\
CURed & B84C4C & Stop conditions and unresolved blockers. \\
CUGreen & 2C8A65 & Successful results. \\
CULight & F3F6FA & Light information surfaces. \\
CUText & 18212E & Body text. \\
CUMuted & 5C6878 & Secondary text. \\
\bottomrule
\end{tabularx}

\subsection{Cover-page standard}

Every cover includes:

\begin{itemize}
    \item CUCII eyebrow in light cyan;
    \item project title in large white type;
    \item document title in light gold;
    \item concise controlled-purpose subtitle;
    \item metadata panel with repository, branches, route, release, technology,
          version, and verification date;
    \item author and repository link.
\end{itemize}

\subsection{Document Control standard}

Page 2 should include:

\begin{itemize}
    \item document name and version;
    \item classification;
    \item repository;
    \item protected branch;
    \item stable branch;
    \item active feature branch when applicable;
    \item baseline or release tag and commit;
    \item target route and key source file;
    \item technical verification date;
    \item Purpose box;
    \item Safety Rule box.
\end{itemize}

\subsection{Section and footer standard}

Section headings use dark navy, bold sans serif, and a cyan rule. Subsections
use dark secondary navy. Subsubsections use blue.

Headers identify the manual and document type. Footers show:

\begin{itemize}
    \item version at left;
    \item page \(x\) of \(y\) centered;
    \item current feature branch or stable release at right.
\end{itemize}

\subsection{Callout boxes}

Use consistent semantic boxes:

\begin{itemize}
    \item Purpose: blue;
    \item Checkpoint: cyan;
    \item Safety Rule: amber;
    \item Do Not Continue Until Resolved: red;
    \item Successful Result: green;
    \item neutral note: light gray-blue;
    \item Authority and Source Priority: gold.
\end{itemize}

Do not use color as the only meaning; each box has a textual title.

\subsection{Prompt and command boxes}

Use:

\begin{itemize}
    \item dark terminal boxes for commands;
    \item blue prompt boxes for copy-and-paste AI instructions;
    \item neutral code boxes for interfaces, schemas, structures, and file
          examples.
\end{itemize}

Prompts should specify:

\begin{enumerate}
    \item repository and branch;
    \item starting commit or tag;
    \item current working-tree expectations;
    \item exact objective;
    \item authorized file scope;
    \item preserve list;
    \item tests and browser review;
    \item strict Git boundaries;
    \item exact final report.
\end{enumerate}

\subsection{Manual section architecture}

A feature implementation manual should normally contain:

\begin{enumerate}
    \item Document Control;
    \item START HERE - new-chat initialization;
    \item governance and verified baseline;
    \item product or mathematical specification;
    \item architecture and file responsibilities;
    \item implementation phases;
    \item automated tests;
    \item accessibility and browser review;
    \item commit, merge, release, and recovery;
    \item appendices with templates and revision record.
\end{enumerate}

\subsection{Source handling}

Clearly distinguish:

\begin{itemize}
    \item what an uploaded source establishes;
    \item what the manual summarizes;
    \item what is a current repository fact;
    \item what is CU interpretation;
    \item what remains unresolved.
\end{itemize}

Do not silently correct research content. Record ambiguities and make explicit
review decisions.

\subsection{Compilation workflow}

\begin{commandbox}{Compile the LaTeX manual}
mkdir -p build

latexmk -pdf \
  -interaction=nonstopmode \
  -halt-on-error \
  -outdir=build \
  Cosmic_Universalism_<Subject>_Implementation_Manual_vX.Y.tex
\end{commandbox}

Run twice automatically through \texttt{latexmk}; inspect the log for
\texttt{Overfull}, \texttt{Underfull}, missing references, and missing files.

\subsection{PDF render and visual QA}

\begin{commandbox}{Render the PDF to PNG pages}
python /home/oai/skills/pdfs/scripts/render_pdf.py \
  build/Cosmic_Universalism_<Subject>_Implementation_Manual_vX.Y.pdf \
  --out_dir build/rendered \
  --dpi 200
\end{commandbox}

Inspect every page for:

\begin{itemize}
    \item clipping or overlap;
    \item broken glyphs;
    \item unreadable code;
    \item orphaned headings;
    \item malformed tables;
    \item header/footer collisions;
    \item incorrect page count;
    \item invalid links or metadata.
\end{itemize}

\subsection{PDF preflight}

\begin{commandbox}{Inspect the final PDF}
python /home/oai/skills/pdfs/scripts/pdf_inspect.py \
  build/Cosmic_Universalism_<Subject>_Implementation_Manual_vX.Y.pdf

python /home/oai/skills/pdfs/scripts/pdf_preflight.py \
  build/Cosmic_Universalism_<Subject>_Implementation_Manual_vX.Y.pdf
\end{commandbox}

\subsection{Manual revision discipline}

Never overwrite the prior published manual without preserving its version.
Record:

\begin{itemize}
    \item version;
    \item date;
    \item source baseline;
    \item sections added or changed;
    \item unresolved issues;
    \item author;
    \item compilation and visual-verification result.
\end{itemize}

% ============================================================
\section{Part 13 - Troubleshooting and Decision Rules}
% ============================================================

\subsection{Wrong branch}

Stop. Report the actual branch and status. Do not switch until the user
authorizes the intended next step, unless the current prompt explicitly
authorized branch creation or switching.

\subsection{Unexpected modified or untracked files}

Do not discard them. Determine whether they are intended work, generated
artifacts, or unrelated files. Ordinary \texttt{git diff} does not show
untracked files; always use:

\begin{commandbox}{List untracked files}
git ls-files --others --exclude-standard | sort
\end{commandbox}

\subsection{Port already in use}

Use the managed Astro status or stop command. Do not use \texttt{--force}
before the normal stop procedure.

\subsection{Tests pass but browser behavior fails}

Treat browser interaction as a real defect. Add a focused regression test where
the architecture permits, but do not install a browser-test dependency without
approval.

\subsection{Manual and repository disagree}

Inspect the live repository and current release tags. Record the discrepancy.
Update the next manual revision; do not silently rewrite history.

\subsection{A source link cannot be opened}

Do not pretend access. Request paste or upload, use embedded context only for
what it supports, and distinguish source-derived facts from CU interpretation.

\subsection{Claude does not enter Aurelius}

Start a fresh Claude conversation and paste the generated prompt again. The
fresh conversation may behave differently even when the prompt structure is
the same. Do not claim that repeated generation guarantees a materially new
prompt.

% ============================================================
\section{Appendix A - Current Project File Manifest}
% ============================================================

\begin{longtable}{L{2.75in}L{3.65in}}
\toprule
\textbf{File or directory} & \textbf{Purpose} \\
\midrule
\filepath{website/src/config/site.ts} & Site identity, nav, header action. \\
\filepath{documentation/master-manual/Cosmic\_Universalism\_Website\_Master\_Implementation\_Manual\_v1.0.tex/.pdf} & Preserved July 21 historical master-manual baseline. \\
\filepath{documentation/master-manual/Cosmic\_Universalism\_Website\_Master\_Implementation\_Manual\_v1.1.tex/.pdf} & Current July 23 master-manual source and compiled PDF. \\
\filepath{documentation/cosmic-breath/ledgers/} & Canonical structural ledger and SHA-256 record. \\
\filepath{website/src/pages/cosmic-breath.astro} & Cosmic Breath route and layered page composition. \\
\filepath{website/src/components/CosmicBreathCycleExplorer.astro} & Interactive structural explorer. \\
\filepath{website/src/components/CosmicBreathDiagram.astro} & Shared static diagram and marker. \\
\filepath{website/src/components/CosmicBreathEducation.astro} & LP-1 field guide. \\
\filepath{website/src/components/CosmicBreathTheoryAndMethod.astro} & LP-2 theory and method. \\
\filepath{website/src/components/CosmicBreathEmpiricalScience.astro} & LP-3 empirical context and sources. \\
\filepath{website/src/components/CosmicBreathPageGuide.astro} & Unified page guide and progressive Back to top. \\
\filepath{website/src/data/cosmic-breath/} & Structural/content authorities, manifest, and 19-source registry. \\
\filepath{website/src/lib/cosmicBreath/} & Runtime projections, engines, atmosphere, quality, and 13 focused tests. \\
\filepath{website/src/pages/cu-time.astro} & CU-Time route. \\
\filepath{website/src/components/CUTimeConverter.astro} & Converter interface. \\
\filepath{website/src/lib/cuTime/} & CU-Time engine and tests. \\
\filepath{website/src/pages/cu-intelligence.astro} & CUCII page and client behavior. \\
\filepath{website/src/components/cucii/} & CUCII presentation components. \\
\filepath{website/src/data/cucii-conversations.ts} & Six Living Conversation records. \\
\filepath{website/src/data/cucii-platforms.ts} & Eleven external AI destinations. \\
\filepath{website/src/data/cucii-prompt-presets.ts} & Proven configurations and menu/purpose presets. \\
\filepath{website/src/data/cucii-sources.ts} & Approved source manifest. \\
\filepath{website/src/data/cucii-working-context.ts} & Embedded CU and Aurelius continuity. \\
\filepath{website/src/lib/cucii/prompt-builder.ts} & Prompt assembler. \\
\filepath{website/src/lib/cucii/validation.ts} & Validation and normalization. \\
\filepath{website/src/lib/cucii/export.ts} & TXT and Markdown exports. \\
\filepath{website/src/lib/cucii/types.ts} & Typed contracts. \\
\filepath{website/src/lib/cucii/__fixtures__/} & Immutable prompt controls. \\
\filepath{website/src/lib/cucii/__tests__/} & CUCII test suite. \\
\bottomrule
\end{longtable}

% ============================================================
\section{Appendix B - Feature Planning Worksheet}
% ============================================================

\begin{longtable}{L{1.8in}L{4.6in}}
\toprule
\textbf{Field} & \textbf{Required entry} \\
\midrule
Feature name & Clear public and technical name. \\
User outcome & What the visitor can do after completion. \\
Baseline & Stable release tag and commit. \\
Feature branch & New branch created from the baseline. \\
Routes & Existing and new routes affected. \\
Authorized files & Exact paths or narrow directories. \\
Protected behavior & Mathematics, URLs, fixtures, tests, and UI behavior that must remain. \\
Data model & Typed registry or interface changes. \\
Accessibility & Keyboard, focus, labels, status, disclosures, and mobile behavior. \\
Tests & Focused and full regression requirements. \\
Browser matrix & 1440, 1024, 768, 390, and 320px. \\
Commit & Proposed commit title and description. \\
Release & Proposed version and annotated tag. \\
Recovery & Previous stable tag. \\
\bottomrule
\end{longtable}

% ============================================================
\section{Appendix C - Final Acceptance Prompt Template}
% ============================================================

\begin{promptbox}{Final Release-Candidate Acceptance}
Repository:
/Users/cosmos/Documents/GitHub/CosmicUniversalismStatement

Required branch:
feature/<feature-name>

Required starting HEAD:
<baseline commit>

The working tree intentionally contains the complete uncommitted
feature.

Perform one final release-candidate review.

Verify:

1. Intended behavior
2. Protected behavior
3. Desktop, tablet, 390px, and 320px layouts
4. Keyboard navigation and visible focus
5. Validation and result-state behavior
6. Reset, copy, download, and external-link behavior
7. Browser console
8. Exact control fixtures and hashes
9. No unrelated file changes

Only correct reproducible defects.
Keep every correction small and add a focused regression test.

Run:

npm test
npm run build
git diff --check
git status --short --branch
git diff --stat
git ls-files --others --exclude-standard | sort

Do not reset, clean, stash, discard, pull, switch branches, commit,
push, merge, tag, release, or deploy.

Return only:

1. Branch and HEAD
2. Files reviewed
3. Confirmed defects
4. Corrections
5. Browser findings
6. Accessibility findings
7. Test count
8. Build result
9. Fixture verification
10. Git status
11. Confirmation that no Git or deployment action occurred
\end{promptbox}

% ============================================================
\section{Appendix D - Release Command Sequence}
% ============================================================

\begin{commandbox}{Feature commit and push sequence}
git branch --show-current
git status --short --branch
git diff --check
git ls-files --others --exclude-standard | sort

git add <exact paths>

git status --short
git diff --cached --name-only
git diff --cached --stat
git diff --cached --check

git commit -m "Add <feature>" -m "<description>"
git push -u origin feature/<feature-name>
\end{commandbox}

\begin{commandbox}{Merge and release sequence}
git switch website
git pull --ff-only origin website

git merge --no-ff feature/<feature-name> \
  -m "Merge <feature>"

cd website
npm test
npm run build
cd ..

git push origin website

git tag -a website-vX.Y-<approved-name> <approved-commit> \
  -m "<approved annotation>"

git push origin refs/tags/website-vX.Y-<approved-name>
\end{commandbox}

% ============================================================
\section{Appendix E - Manual Source Package Template}
% ============================================================

A future manual project should create:

\begin{codebox}{Recommended documentation package}
<Manual_Name>_vX.Y.tex
<Manual_Name>_vX.Y.pdf
CURRENT_STATE_HANDOFF.md
SOURCE_MANIFEST.md
<feature-control-specimen>
\end{codebox}

When only five ChatGPT Project uploads are available, the Project Instructions
and first-chat bootstrap remain text in settings/chat and do not consume file
slots.

\subsection*{SOURCE\_MANIFEST.md template}

\begin{codebox}{Source manifest}
# Source Manifest

Manual:
<manual filename and version>

Repository baseline:
<tag and commit>

Primary source files:
- <filename>
- <filename>

Control specimens:
- <path, byte length, SHA-256>

Historical manuals consulted:
- <filename and version>

Current facts verified from:
- git branch --show-current
- git rev-parse --short HEAD
- git status --short --branch
- npm test
- npm run build

Unresolved source conflicts:
- <none or list>

Compilation:
- pdfLaTeX / latexmk
- render verified at 200 DPI
\end{codebox}

% ============================================================
\section{Appendix F - Canonical LaTeX Skeleton}
% ============================================================

\begin{codebox}{Minimal future manual skeleton}
\documentclass[11pt,letterpaper]{article}

\usepackage[margin=0.78in,headheight=16pt]{geometry}
\usepackage[T1]{fontenc}
\usepackage[utf8]{inputenc}
\usepackage{lmodern}
\usepackage[scaled=0.94]{helvet}
\usepackage{microtype}
\usepackage{parskip}
\usepackage{setspace}
\usepackage{enumitem}
\usepackage{booktabs}
\usepackage{tabularx}
\usepackage{xcolor}
\usepackage{graphicx}
\usepackage{fancyhdr}
\usepackage{lastpage}
\usepackage{titlesec}
\usepackage{hyperref}
\usepackage{xurl}
\usepackage{bookmark}
\usepackage{listings}
\usepackage[most]{tcolorbox}

\renewcommand{\familydefault}{\sfdefault}
\setstretch{1.08}
\setlength{\parindent}{0pt}
\setlength{\parskip}{0.65em}

% Copy the canonical CU color definitions, listing styles,
% tcolorbox definitions, title formatting, cover, and Document
% Control structure from this master .tex source.

\begin{document}

% Cover page
% Document Control
% START HERE
% Governance
% Specification
% Architecture
% Tests
% Release
% Appendices
% Revision record

\end{document}
\end{codebox}

\begin{authoringbox}
Do not recreate the house style from memory. Copy the preamble and structural
patterns from the delivered master \texttt{.tex} file, then change only the
metadata, subject-specific content, branch/release details, and required
appendices.
\end{authoringbox}

% ============================================================
\section{Appendix G - Source Provenance}
% ============================================================

This manual consolidates and updates the following project sources:

\begin{itemize}
    \item \filepath{Cosmic\_Universalism\_Website\_Implementation\_Manual\_v1.1.pdf/.tex};
    \item \filepath{Cosmic\_Breath\_Time\_Converter\_Implementation\_Manual\_v1.0.pdf/.tex};
    \item \filepath{CUCII\_AI\_Exploration\_and\_Prompt\_Studio\_Implementation\_Manual\_v1.0.pdf/.tex};
    \item \filepath{Cosmic\_Universalism\_Cosmic\_Breath\_Interactive\_Cycle\_Explorer\_Implementation\_Manual\_v1.1.pdf/.tex};
    \item \filepath{cosmic-universalism-website-blueprint.md};
    \item \filepath{cucii-homepage-inspection.txt};
    \item \filepath{cosmic\_breath\_time\_converter\_v3\_0\_1.html};
    \item \filepath{Cosmic\_Breath\_Calculation.md};
    \item \filepath{Cosmic\_Breathing\_Cycle.md};
    \item \filepath{Time\_Calculation.md};
    \item \filepath{CUCII\_Website\_v1.2\_New\_Account\_Handoff.md};
    \item \filepath{Aurelius\_Claude\_Working\_Prompt\_v1.0.txt};
    \item the approved \release{CB-TOM-STRUCTURAL-1.0} and \release{CB-TOM-CONTENT-1.0} authorities and their digest records;
    \item verified Cosmic Breath implementation, QA, branch-comparison, merge, push, and tag outputs through
          \release{website-v1.1-cosmic-breath} at
          \filepath{7ac87167c51aad9ae7a31e17d1b2127fc069e15c}.
\end{itemize}

Earlier manuals are preserved as historical implementation records. This
master manual is the current operational authority as of its verification
date.

% ============================================================
\section{Appendix H - Revision Record}
% ============================================================

\begin{tabularx}{\textwidth}{L{0.7in}L{1.25in}Y}
\toprule
\textbf{Version} & \textbf{Date} & \textbf{Revision} \\
\midrule
1.0 & July 21, 2026 & Initial unified master manual. Consolidates the website,
CU-Time, and CUCII manuals; records the v1.3 Aurelius release; defines current
ChatGPT/Codex project governance; and establishes the canonical LaTeX house
style and project-file templates for future manuals. \\
1.1 & July 23, 2026 & Adds the completed Cosmic Breath Cycle Explorer; sealed
structural and companion-content authorities; browser-safe projections;
guarded transitions; LP-1, LP-2, and LP-3 learning layers; unified navigation;
19 scientific sources; 251-test verification; merge and tag records; canonical
digests; and the current website integration baseline. \\
\bottomrule
\end{tabularx}

\vfill
\begin{center}
\textcolor{CUMuted}{End of Cosmic Universalism Website Master Implementation and Operations Manual}
\end{center}

\end{document}
